\subsection{How Do We Break a Vicious Cycle?}

If you want to break a vicious cycle, do not start by arguing with the cycle.
Start by changing the conditions that keep re-creating it. In this framework,
the cycle persists because the current behavioural frame keeps winning against
the target state whenever the barrier is too high, the decision window is
missed, or volitional energy is spent on the wrong move. The practical answer
is to lower the barrier to the next target state, act inside the available
decision window, and install one small change that makes the next repetition
cheaper.

\begin{enumerate}[nosep]
  \item Name the cycle in concrete terms. What repeats, when does it repeat,
        and what state keeps winning?
  \item Identify the target state for the next loop. Do not aim at a perfect
        life. Aim at the next stabilizing state that can actually hold.
  \item Find the dominant barrier. Is the problem unclear first action,
        exhausted volitional energy, a hostile environment, or delay until the
        decision window has already closed?
  \item Reduce one barrier variable before the next attempt. Move the object,
        prewrite the first line, change the room, or shrink the entry action so
        the target state becomes easier to enter.
  \item Act inside the decision window with one visible step. The point is not
        to solve the entire cycle in one move, but to prevent the old frame from
        rebuilding itself.
  \item Repeat the same low-friction setup until the new pattern has history.
        A cycle weakens when the alternative state becomes easier to start than
        the old one.
\end{enumerate}

\begin{remark}
A student who keeps postponing study may treat the problem as a lack of moral
strength. Mechanically, the more useful diagnosis is often that the target
state is too vague, the barrier is too high, and the decision window is being
used for thinking instead of entering. The fix is usually small: open the file
before the slump, define the first task in one sentence, and begin while the
window is still open.
\end{remark}
